\documentclass[10pt, aspectratio=169]{beamer}
\usetheme{Madrid}
\definecolor{myGreen}{RGB}{46,110,64}
\definecolor{myLightGreen}{RGB}{224,240,224}
\usecolortheme[named=myGreen]{structure}
\setbeamertemplate{footline}{%
  \leavevmode%
  \hbox{\begin{beamercolorbox}[wd=.5\paperwidth,ht=2.25ex,dp=1ex,center]{author in head/foot}
    \insertshortauthor
  \end{beamercolorbox}%
  \begin{beamercolorbox}[wd=.5\paperwidth,ht=2.25ex,dp=1ex,right]{date in head/foot}
    \insertframenumber{} / \inserttotalframenumber\hspace*{2ex}
  \end{beamercolorbox}}%
}
\usepackage{graphicx}
\usepackage{amsmath}
\usepackage{tcolorbox}
\tcbuselibrary{skins}

% Reusable box for highlighted key numbers
\newtcbox{\keystat}{on line, colback=myGreen!12, colframe=myGreen, boxrule=0.5pt,
  arc=2pt, boxsep=1pt, left=4pt, right=4pt, top=2pt, bottom=2pt}

\title[NDVI vs. NBR]{NDVI vs. NBR}
\subtitle{Assessing Post-Fire Vegetation Recovery Monitoring at Hermits Peak}
\author[F. Hou]{Francine Hou}
\institute[ICR]{Institute for Computing in Research}
\date{July 2026}

\begin{document}

% ============================================================
\frame{\titlepage}

% ============================================================% SLIDE 1: Context & Disaster Overview
% -----------------------------------------------------------------------------
\begin{frame}{Context: The Hermits Peak / Calf Canyon Fire}
    \begin{block}{Background Context}
        In April 2022, two prescribed burns merged to form the \textbf{Hermits Peak / Calf Canyon Fire}, which consumed over \textbf{341,000 acres} across northern New Mexico to become the largest wildfire in the state's recorded history.
    \end{block}

    \vspace{0.5em}

    \begin{block}{Economic \& Environmental Impact}
        Ultimately, the fire cost the displacement of thousands of residents and over 5 billion dollars in recovery efforts.
        Note: High-severity fires can trigger long term shifts of forest ecosystems into shrubs and grasslands, which are prone to experiencing more frequent fires.
    \end{block}

    
\end{frame}

% --------------============================================================
\begin{frame}{The Research Gap}

\begin{itemize}
    \item \textbf{The Santa Fe Forest Reserve} relies on \textbf{NDVI alone} for post-fire canopy monitoring and vegetation assessments at Hermits Peak.
\end{itemize}

\vspace{0.5em}

\begin{alertblock}{The Missing Comparative Knowledge}
\begin{itemize}
    \item \textbf{Limited multi-index validation:} NDVI is vulnerable to soil background interference and early  grass/shrub regrowth, which can inflate its readings after a fire.
    \item \textbf{No direct comparative study:} no empirical comparison between NDVI and NBR has been conducted specifically for the Hermits Peak--Calf Canyon burn scar.
\end{itemize}
\end{alertblock}
\end{frame}

% ============================================================


\begin{frame}{NDVI vs NBR}
\begin{itemize}
    \item \textbf{NDVI (Normalized Difference Vegetation Index)} Used to assess vegetation health and density.  Near-infrared (NIR) reflectance and red light absorption. 
    \item \textbf{NBR (Normalized Burn Ratio)} Used to highlight burnt areas from active or recent wildfires. NIR reflectance and shortwave infrared (SWIR) absorption (more sensitive to \alert{moisture content and char}).

    
\end{itemize}
\end{frame}
\begin{frame}{Research Hypotheses}
\begin{block}{Main Hypothesis ($H_1$)}
{NDVI alone will overstate post-fire recovery} in the Hermits Peak burn scar as NDVI values are consistently higher than NBR. 
\end{block}

\vspace{0.8em}

\begin{exampleblock}{Null / Alternative Hypothesis ($H_0$)}
{NDVI provides sufficient accuracy} for long-term post-fire monitoring on its own, rendering an additional index like NBR statistically redundant for Forest Service management purposes.
\end{exampleblock}
\end{frame}

% ============================================================
% ============================================================
\begin{frame}{Data Sources}
\begin{itemize}
    \item Imagery source: \textbf{Landsat 8/9 Collection 2, Level-2} surface reflectance (USGS EarthExplorer)
    \item Location: \keystat{Path 33, Row 35} -- covering the Hermits Peak burn scar
    \item Time series: one \textbf{May} scene per year, \keystat{2021 -- 2026}, giving a pre-fire baseline plus five recovery years
    \item Quality filter: cloud cover \keystat{$<$15\%} for every scene used
    \item Spatial extent: all rasters clipped to the official burn perimeter from the \textbf{BAER} Imagery Support Program before analysis
\end{itemize}
\end{frame}

\begin{frame}{NDVI NBR Calculation}
\small
Red, near-infrared, and shortwave infrared bands were extracted from each scene and combined using the following formulas:

\begin{center}
\begin{tcolorbox}[colback=myLightGreen, colframe=myGreen, arc=3pt, boxrule=0.8pt, width=0.7\textwidth]
\centering
\vspace{1pt}
$\displaystyle \text{NDVI} = \frac{\text{NIR} - \text{Red}}{\text{NIR} + \text{Red}}$
\vspace{3pt}
\end{tcolorbox}

\vspace{4pt}

\begin{tcolorbox}[colback=myLightGreen, colframe=myGreen, arc=3pt, boxrule=0.8pt, width=0.7\textwidth]
\centering
\vspace{1pt}
$\displaystyle \text{NBR} = \frac{\text{NIR} - \text{SWIR2}}{\text{NIR} + \text{SWIR2}}$
\vspace{3pt}
\end{tcolorbox}
\end{center}

\begin{itemize}
    \item Calculated in \textbf{Python} (\texttt{rasterio}, \texttt{numpy}) from downloaded Landsat bands
    \item Symbolized and spatially processed in \textbf{QGIS}
    \item 500 random points sampled within the burn boundary
\end{itemize}
\end{frame}

% ============================================================
\begin{frame}{NDVI, 2021--2026}
\centering
\includegraphics[height=0.32\textheight]{2021ndvi.jpeg}
\includegraphics[height=0.32\textheight]{2022ndvi.jpeg}
\includegraphics[height=0.32\textheight]{2023ndvi.jpeg}\\[2pt]
\includegraphics[height=0.34\textheight]{2024ndvi.jpeg}
\includegraphics[height=0.34\textheight]{2025ndvi.jpg}
\includegraphics[height=0.34\textheight]{2026ndvi.jpg}

\vspace{4pt}
\footnotesize 2021 (pre-fire) through 2026, left to right, top to bottom.
\end{frame}

\begin{frame}{NBR, 2021--2026}
\centering
\includegraphics[height=0.38\textheight]{2021nbr.jpeg}
\includegraphics[height=0.38\textheight]{2022nbr.jpeg}
\includegraphics[height=0.38\textheight]{2023nbr.jpeg}\\[2pt]
\includegraphics[height=0.38\textheight]{2024nbr.jpeg}
\includegraphics[height=0.38\textheight]{2025nbr.jpeg}
\includegraphics[height=0.38\textheight]{2026nbr.jpeg}
\end{frame}


% ============================================================

\begin{frame}{Comparison and Analysis}
\centering
\includegraphics[height=0.75\textheight]{555.png}

\end{frame}

\begin{frame}{Comparison and Analysis: dNDVI and dNBR}
\small
Change relative to a pre-fire baseline (2021).

\vspace{8pt}
\begin{columns}[T]
    \begin{column}{0.48\textwidth}
        \begin{tcolorbox}[colback=myLightGreen, colframe=myGreen, arc=3pt,
            boxrule=0.8pt, left=6pt, right=6pt, top=4pt, bottom=4pt]
        \centering
        \textbf{dNDVI}\\[4pt]
        $\text{dNDVI} = \text{NDVI}_{\text{year}} - \text{NDVI}_{2021}$
        \end{tcolorbox}
    \end{column}

    \begin{column}{0.48\textwidth}
        \begin{tcolorbox}[colback=myLightGreen, colframe=myGreen, arc=3pt,
            boxrule=0.8pt, left=6pt, right=6pt, top=4pt, bottom=4pt]
        \centering
        \textbf{dNBR}\\[4pt]
        $\text{dNBR} = \text{NBR}_{\text{year}} - \text{NBR}_{2021}$
        \end{tcolorbox}
    \end{column}
\end{columns}

\vspace{10pt}
\begin{itemize}
    \item Calculated \textbf{per sample point}, for every post-fire year (2022--2026)
    \item A value near \textbf{0} means that point has returned to its pre-fire condition
    \item Differencing removes that baseline noise and highlights {fire-driven change}, which a single NBR/NDVI image cannot do effectively on its own.
\end{itemize}

\end{frame}



% ============================================================
\begin{frame}{Comparison and Analysis}
\centering
\includegraphics[height=0.75\textheight]{333.png}

\end{frame}

\begin{frame}{Comparison and Analysis}
\centering
\includegraphics[height=0.78\textheight]{444.png}

\end{frame}

\begin{frame}{Comparison and Analysis}
\small
\begin{itemize}
    \item Lowest NDVI--NBR correlation of the entire study period corresponds to the \textbf{first full growing season} after containment
    \item Exactly when fast-colonizing grasses and shrubs are expected to green up, well ahead of any real canopy return
\end{itemize}

\centering
\includegraphics[height=0.55\textheight]{222.png}

\footnotesize Notice the wider spread around the trend line in the 2023 panel compared to other years.
\end{frame}

% ============================================================
\begin{frame}{Comparison and Analysis}
\small
\begin{itemize}
    \item Each of the 500 sample points was compared to its own {2021 baseline} value
    \item Classified as {fully recovered} (80-100), {partially recovered} (50-80), or {not recovered} (0-50) 
\end{itemize}

\centering
\includegraphics[height=0.58\textheight]{111.png}
\end{frame}

\begin{frame}{Comparison and Analysis}
\begin{itemize}
    \item \textbf{By NDVI:} \keystat{62\% fully recovered} \quad \keystat{$<$2\% not recovered}
    \item \textbf{By NBR:} \keystat{31\% fully recovered} \quad \keystat{52\% not recovered}
    \item By NDVI's measure, the landscape has largely healed. By NBR's measure, most of it has not.
\end{itemize}
\end{frame}

% ============================================================
\begin{frame}{Implications}
\begin{itemize}
    \item NDVI and NBR are highly correalated, but differences in recovery status make them not interchangeable. 
    \item The gap between them is largest exactly when it matters most: the first years after containment, when management decisions about restoration are often made.
    \item A monitoring approach built on NDVI alone risks declaring a landscape recovered well before the forest structure has actually returned.
    \item Pairing NDVI with a burn-sensitive index like NBR would give the forest reserve a more complete picture of recovery.
\end{itemize}
\end{frame}

\begin{frame}{Limitations and Future Research}
\begin{itemize}
    \item Only \textbf{six} annual time points -- one scene per year
    \item No field validation of recovery classifications
    \item 500 random points, \textbf{not} stratified by burn severity zone
    \item No canopy height or structural data (e.g., LiDAR)
    \item Single month (May) per year -- seasonal effects not fully explored
\end{itemize}

\vspace{6pt}
\textbf{Next steps:} field validation, GEDI LiDAR, canopy cover metrics.
\end{frame}

% ============================================================
\begin{frame}{Conclusion}
\begin{itemize}
    \item NDVI and NBR agree on the \textbf{direction} of recovery and demonstrate relatviely high correlation.
    \item NDVI consistently reads more optimistically than NBR.
    \item For future monitoring of the Hermits Peak/Calf Canyon burn scar and possibly other fires, it is recommended to include additional indices and metrics.  
\end{itemize}

\vspace{16pt}
\centering
\Large Thank you.
\end{frame}

\end{document}